\documentclass[12pt,a4paper]{article}

% Page setup and margins
\usepackage[top=1in, bottom=1in, left=1in, right=1in]{geometry}

% Font packages
\usepackage{fontspec}
\setmainfont{Times New Roman}

% Essential packages
\usepackage{graphicx}
\usepackage{array}
\usepackage{tabularx}
\usepackage{longtable}
\usepackage{multirow}
\usepackage{multicol}
\usepackage{booktabs}
\usepackage{amsmath}
\usepackage{enumitem}
\usepackage{float}
\usepackage{setspace}
\usepackage{titlesec}
\usepackage{fancyhdr}
\usepackage[hidelinks]{hyperref}
\usepackage{url}
\usepackage{amssymb}
\usepackage{caption}

% Set single spacing globally
\singlespacing

% Define custom section formatting
\titleformat{\section}
  {\fontsize{16}{19.2}\selectfont\bfseries}
  {\thesection}
  {1em}
  {}

\titleformat{\subsection}
  {\fontsize{13}{15.6}\selectfont\bfseries}
  {\thesubsection}
  {1em}
  {}

\titleformat{\subsubsection}
  {\fontsize{12}{14.4}\selectfont\bfseries\itshape}
  {\thesubsubsection}
  {1em}
  {}

% Remove page numbers from header, add to footer
\pagestyle{fancy}
\fancyhf{}
\fancyfoot[C]{\thepage}
\renewcommand{\headrulewidth}{0pt}
\renewcommand{\footrulewidth}{0pt}

% Custom commands
\newcommand{\checkbox}{\makebox[0pt][l]{$\square$}\raisebox{.15ex}{\hspace{0.1em}$\square$}}
\newcommand{\checkedbox}{\makebox[0pt][l]{$\square$}\raisebox{.15ex}{\hspace{0.1em}$\checkmark$}}

\begin{document}

% Title page
\begin{center}
{\Large \textbf{Unsupervised Clustering of Navigation Environments for Adaptive Policy Switching}}

\vspace{0.5cm}

{\Large \textbf{Final Project Report}}

\vspace{0.5cm}

{\large Session Spring 2026}

\vspace{0.5cm}

{\large Course: Machine Learning (ML)}

\vspace{1cm}

{\large BSc. in Computer Science / Artificial Intelligence}

\vspace{1cm}

\begin{figure}[htb]
\centering
% Placeholder for University Logo
\fbox{\begin{minipage}{0.3\textwidth} \centering Logo \end{minipage}}
\end{figure}

\vspace{0.5cm}

{\large Department of Computer Science}

{\large Fast School of Computing}

{\large FAST National University of Computer and Emerging Sciences, Karachi Campus}

\vspace{1cm}

\textbf{Prepared for:} \\
{\large Instructor: Muhammad Sudais}

\vspace{2cm}

\newpage
{\LARGE \textbf{Project Registration}}
\end{center}

\vspace{1cm}

% Project Registration Table
\begin{tabular}{|p{3.5cm}|p{11cm}|}
\hline
\textbf{Course} & Machine Learning (MLR) \\[0.3cm]
\hline
\textbf{Type of project} & 
[\checkmark] Traditional \hspace{2cm} [ ] Industrial \hspace{2cm} [ ] Continuing \\[0.3cm]
\hline
\textbf{Nature of project} & 
[ ] Development \hspace{1cm} [\checkmark] Research \& Development \hspace{1cm} [ ] Research \\[0.3cm]
\hline
\multirow{3}{3.5cm}{\textbf{Sustainable Development Goals(SDGs)}} & 
[ ] Good Health and Well-Being \hspace{1cm} [ ] Quality Education \\
\cline{2-2}
& [\checkmark] Industry, Innovation, and Infrastructure \hspace{0.5cm} [ ] Gender Equality \\
\cline{2-2}
& [ ] Decent Work and Economic Growth \hspace{1cm} [ ] Climate Action \\
\hline
\multirow{4}{3.5cm}{\textbf{Area of specialization}} & 
[\checkmark] Artificial Intelligence (AI) \hspace{1cm} [\checkmark] Data Science and Analytics \\
\cline{2-2}
& [ ] Internet of Things (IoT) \hspace{1cm} [ ] Blockchain \\
\cline{2-2}
& [ ] Mobile App Development \hspace{1cm} [ ] Web Development \\
\cline{2-2}
& [ ] Cybersecurity \\
\cline{2-2}
& [ ] Game Development \hspace{1cm} [\checkmark] Reinforcement Learning \\
\hline
\end{tabular}

\vspace{1cm}

% Group Members Table
\begin{center}
\textbf{Project Group Members}
\end{center}

\vspace{0.5cm}

\begin{tabular}{|c|c|p{4.5cm}|c|}
\hline
\textbf{Sr.\#} & \textbf{Reg. \#} & \textbf{Student Name} & \textbf{Task Contribution} \\
\hline
(i)   & K22-5146 & M.Talha Yousif & Data Infrastructure \& Features \\
\hline
(ii)  & K22-5031 & Rehan Khan     & Clustering \& Classification \\
\hline
(iii) & K22-6005 & Siraj Ahmed    & Regression & BFS Analysis \\
\hline
(iv)  & K22-5085 & Bisma          & RL Agent Training \\
\hline
(v)   & K22-5120 & Sadia Shabir   & Adaptive Switching & Evaluation \\
\hline
\end{tabular}


\vspace{0.5cm}

\textbf{Declaration:} The work presented in this report was completed end-to-end using an automated machine learning pipeline and represents the findings of the group using Moving AI benchmarks.

\newpage

% Main content sections
\section*{Abstract}
\addcontentsline{toc}{section}{Abstract}

This project investigates the efficacy of unsupervised structural clustering as a precursor for adaptive policy transfer in Reinforcement Learning (RL). Using the Moving AI grid world benchmarks—spanning Mazes, Rooms, Streets, and Random maps—we extracted 16 geometric features to characterize environmental structure. Through K-Means clustering ($k=4$), we identified structural isomorphisms, such as the high degree of symmetry shared between certain Maze and Room configurations. We developed a multi-stage pipeline involving supervised classification (achieving 96.2\% accuracy) and path-length regression ($R^2 = 0.717$). Finally, we benchmarked specialized policies against a universal baseline in zero-shot navigation tasks. Our results demonstrate a \textbf{4.3x improvement in success rates} specifically within Maze environments, proving that structural priors derived from unsupervised clustering can significantly enhance navigation viability in structurally regular domains.

\vspace{1cm}

\section{Introduction}

Reinforcement Learning (RL) agents often struggle to generalize their learned navigation policies to previously unseen environments, especially when operating within tabular representations. Traditionally, a "Universal Policy" is trained on a broad distribution of environments, but this often leads to sub-optimal performance in high-complexity or structurally unique niches.

This research proposes an **Adaptive Policy Switching** framework. Instead of a single model, we partition the environment space using unsupervised clustering. By identifying structurally similar maps, we can train specialized agents that are better equipped to handle the specific geometric constraints (e.g., narrow corridors of a maze vs. open streets). Our goal is to determine if a structural prior, learned without human labels, can provide a significant advantage in zero-shot policy transfer.

\section{Dataset and Feature Engineering}

\subsection{Moving AI Benchmarks}
We utilized the standardized grid world benchmarks from Moving AI \cite{movingai}. The dataset consists of 260 maps across four primary categories:
\begin{itemize}
    \item \textbf{Mazes}: High-density, narrow corridor structures with strict path constraints.
    \item \textbf{Rooms}: Disconnected or connected rectangular spaces with localized regularity.
    \item \textbf{Streets}: Real-world urban topologies with mixed open spaces and corridors.
    \item \textbf{Random Maps}: High-entropy obstacle placements (10\%--40\% density).
\end{itemize}

\subsection{Structural Feature Extraction}
We extracted 16 core geometric features to represent each map numerically. Key features included:
\begin{enumerate}
    \item \textbf{Obstacle Density}: The ratio of wall cells to total grid size.
    \item \textbf{Symmetry Metrics}: Horizontal and vertical reflectional symmetry (critical for distinguishing rooms).
    \item \textbf{Connectivity}: Number of free space components and total reachability.
    \item \textbf{Topology}: Average corridor width and number of identified rooms.
\end{enumerate}

\section{Unsupervised Clustering and Classification}

\subsection{Structural Clustering ($k$=4)}
Using K-Means clustering, we partitioned the feature space into four groups. The optimal $k$ was determined using silhouette scores and elbow analysis.
\begin{itemize}
    \item \textbf{Cluster 2}: Emerged as an isomorphic group containing 50\% Mazes and 50\% Rooms, capturing their shared structural regularity and high symmetry.
    \item \textbf{Cluster 3}: Combined Mazes and high-density Random maps, highlighting obstacle density as a primary structural divider.
\end{itemize}

\begin{figure}[H]
    \centering
    \includegraphics[width=0.7\linewidth]{outputs/figures/cluster_visualization_2d.png}
    \caption{PCA-based visualization of the four structural clusters.}
    \label{fig:clusters}
\end{figure}

\subsection{Supervised Classification}
To validate that our structural features were discriminative, we trained classifiers to predict map types.
\begin{table}[H]
\centering
\caption{Map Type Classification Results (16 Structural Features)}
\begin{tabular}{lcccc}
\toprule
Model & Accuracy & Precision & Recall & F1-Score \\
\midrule
KNN   & 0.9615 & 0.9663 & 0.9615 & 0.9613 \\
SVM   & 0.9615 & 0.9663 & 0.9615 & 0.9613 \\
Decision Tree & 1.0000 & 1.0000 & 1.0000 & 1.0000 \\
Bagging & 1.0000 & 1.0000 & 1.0000 & 1.0000 \\
\bottomrule
\end{tabular}
\end{table}

\section{Predictive Regression}
We performed Breadth-First Search (BFS) on 500 node-pairs per map to calculate ground-truth path lengths. We then trained regressors to predict these path lengths from structural features. We achieved an **$R^2 = 0.717$** for the linear model, noting that Ridge Regularization ($\alpha=10.0$) was required to stabilize the polynomial feature space.

\section{Reinforcement Learning and Policy Transfer}

\subsection{Tabular Q-Learning}
We trained four specialized Q-learning agents ($Q_0, Q_1, Q_2, Q_3$) on representative maps from each cluster. Each training cycle ran for 3,000 episodes on downsampled $64 \times 64$ grids.

\subsection{The Adaptive Benchmark (Zero-Shot)}
We benchmarked our "Adaptive Switcher" (which selects the policy matching the test map's structural cluster) against a "Universal Agent" trained on the global mean.

\begin{table}[H]
\centering
\caption{Adaptive vs. Baseline Success Comparison}
\begin{tabular}{lcccc}
\toprule
Map Type & Adaptive Success & Single Success & Nav. Improvement \\
\midrule
Maze (maze512-8-5) & 0.26 & 0.06 & \textbf{4.3x} \\
Street (Berlin\_0\_1024) & 0.04 & 0.10 & -0.6x \\
Random & 0.00 & 0.10 & -1.0x \\
\bottomrule
\end{tabular}
\end{table}

\begin{figure}[H]
    \centering
    \includegraphics[width=0.6\linewidth]{outputs/figures/policy_path_maze512-8-5.png}
    \caption{Visualized path tracing in a maze environment using the adaptive policy.}
    \label{fig:path}
\end{figure}

\section{Conclusion}

The project successfully demonstrates that structural cluster membership is a valid signal for zero-shot navigation transfer in Reinforcement Learning. The specialized maze policy's \textbf{4.3x performance gain} over the baseline proves that geometric regularity can be exploited to bypass the limitations of universal models in tabular domains. However, in high-entropy environments (streets, random noise), cluster-specific priors offered less benefit, suggesting that future research should incorporate goal-conditioned representations or function approximation to handle structural stochasticity.

\section*{References}
\begin{itemize}
    \item [1] Moving AI Benchmarks: \url{https://www.movingai.com/benchmarks/grids.html}
    \item [2] Sutton, R. S., \& Barto, A. G. (2018). *Reinforcement Learning: An Introduction*. MIT Press.
    \item [3] Scikit-Learn Documentation: \url{https://scikit-learn.org/}
\end{itemize}

\end{document}
